% =============================================================================
% APPENDIX S: PREDICT-FALSIFIABLE: Ten Testable Hypotheses
% =============================================================================
% Module: EBF_S_PREDICT
% Category: PREDICT
% Version: 1.0 (January 2026)
% Status: Draft
% Language: English
%
% This appendix derives specific, falsifiable predictions from the EBF framework
% to address the Popperian challenge of scientific testability.
%
% =============================================================================

% =============================================================================
% CROSS-REFERENCE STRUCTURE
% =============================================================================
%
% CHAPTER LINKAGE (Required):
% - Primary Chapter: Chapter 12: Empirical Validation
% - Secondary Chapters: Chapter 3 (Complementarity), Chapter 5 (Context)
%
% This appendix provides foundation for empirical testing:
%
% DEPENDENCIES (what this appendix REQUIRES):
% - Appendix HOW (CORE-HOW): Complementarity structure $C$
% - Appendix CTW (CORE-WHEN): Context dimensions $\Psi$
% - Appendix WAT (CORE-WHAT): Utility dimensions (FEPSDE)
% - Appendix CAM (Metatheory): Framework philosophy
%
% DEPENDENTS (what REQUIRES this appendix):
% - Appendix THL1 (Evaluation Protocol): Uses predictions as test cases
% - Appendix EPS (Epistemics of Deviation): Uses failed predictions for learning
% - Chapter 12: References predictions for validation design
%
% =============================================================================

% -----------------------------------------------------------------------------
% CHAPTER-APPENDIX LINKAGE BOX
% -----------------------------------------------------------------------------

\begin{tcolorbox}[colback=purple!5!white, colframe=purple!75!black,
    title=Chapter Linkage]

\textbf{Primary Chapter:} Chapter 12: Empirical Validation
\begin{itemize}[nosep]
\item Section 12.2: Testable Hypotheses $\leftrightarrow$ This Appendix Section S.2-S.11
\item Section 12.4: Falsification Criteria $\leftrightarrow$ This Appendix (each prediction)
\end{itemize}

\textbf{Secondary Chapters:}
\begin{itemize}[nosep]
\item Chapter 3: Complementarity Theory --- references predictions 1-3, 5-6
\item Chapter 5: Context Modulation --- references predictions 4, 7-8
\item Chapter 7: Identity and Welfare --- references prediction 9
\end{itemize}

\textbf{Reading Guide:}
\begin{itemize}[nosep]
\item For theoretical foundation: Read Chapters 3 and 5 first
\item For testing methodology: See Appendix THL1 (Evaluation Protocol)
\item For learning from failures: See Appendix EPS (Epistemics of Deviation)
\end{itemize}

\end{tcolorbox}

% -----------------------------------------------------------------------------
% CHAPTER TITLE + LABEL
% -----------------------------------------------------------------------------

\chapter{PREDICT-FALSIFIABLE: Ten Testable Hypotheses}
\label{app:S-falsifiable}


\begin{tcolorbox}[colback=blue!5!white, colframe=blue!75!black,
    title=Quick Reference: Appendix SUN]

\textbf{Appendix:} SUN (PREDICT)

\textbf{Core Question:} What are the key findings in this domain?

\textbf{Key Concepts:}
\begin{itemize}[nosep]
\item Framework integration
\item Parameter validation
\item Cross-references
\end{itemize}

\textbf{Cross-References:} See related appendices below
\end{tcolorbox}

% -----------------------------------------------------------------------------
% ABSTRACT
% -----------------------------------------------------------------------------

\begin{abstract}
A framework that can explain everything explains nothing. This appendix addresses the Popperian challenge by deriving ten specific, falsifiable predictions from the EBF $(C, \Psi, C^*, K, Q)$ framework. Each prediction specifies the direction and magnitude of effects, identifies what evidence would refute it, and provides concrete empirical tests. Topics include cash transfer effectiveness, management practice transfer, information intervention asymmetry, financial-political crisis lags, recovery convexity, trade war hysteresis, technology-culture divergence, multi-level coherence conflicts, identity-economics non-substitutability, and coherent traps.
\end{abstract}

% -----------------------------------------------------------------------------
% QUICK REFERENCE BOX
% -----------------------------------------------------------------------------

\begin{tcolorbox}[colback=blue!5!white, colframe=blue!75!black,
    title=Quick Reference: Falsifiable Predictions]

\textbf{Core Question:} What specific, testable predictions does EBF generate that could falsify the framework?

\textbf{Key Formula:}
\begin{equation}
Q = f(C^*(\Psi), K) \quad \text{where } C^*(\Psi) = C \cdot \Psi
\end{equation}

\textbf{Key Concepts:}
\begin{itemize}[nosep]
\item Falsifiability: Each prediction specifies what evidence would refute it
\item Context Moderation: $\Psi$ determines which complementarities ($C$) are active
\item Coherence: $K$ measures alignment with effective complementarity structure $C^*$
\end{itemize}

\textbf{Cross-References:} Appendix THL1 (Evaluation Protocol), Appendix CAM (Metatheory), Appendix EPS (Epistemics of Deviation), Appendix FRM (Philosophy of Science Foundation---Fehr Compatibility)
\end{tcolorbox}

% -----------------------------------------------------------------------------
% APPENDIX SCOPE BOX
% -----------------------------------------------------------------------------

\begin{tcolorbox}[
    colback=orange!5!white,
    colframe=orange!75!black,
    title={\textbf{Appendix Scope: Ziel / In-Scope / Out-of-Scope / Constraints}},
    fonttitle=\bfseries
]

\textbf{Ziel (Objective):}
\begin{itemize}[nosep]
    \item The critique
\end{itemize}

\textbf{In-Scope:}
\begin{itemize}[nosep]
    \item The Fundamental Question
    \item Prediction 1: Cash Transfer Effectiveness and Social Capital
    \item Prediction 2: Management Practice Transfer Failure
    \item Prediction 3: Information Intervention Asymmetry
    \item Prediction 4: Financial Crisis Political Lag
\end{itemize}

\textbf{Out-of-Scope (delegiert an andere Appendices):}
\begin{itemize}[nosep]
    \item CORE-HOW $\rightarrow$ Appendix HOW
    \item CORE-WHEN $\rightarrow$ Appendix CTW
    \item CORE-WHAT $\rightarrow$ Appendix WAT
    \item Metatheory $\rightarrow$ Appendix CAM
    \item Epistemics of Deviation $\rightarrow$ Appendix EPS
\end{itemize}

\textbf{Constraints (Anwendungsgrenzen):}
\begin{itemize}[nosep]
    \item [TODO: Define when this appendix does NOT apply]
    \item [TODO: Define required prerequisites]
    \item [TODO: Define known limitations]
\end{itemize}

\textbf{Lieferobjekte (Deliverables):}
\begin{itemize}[nosep]
    \item 2 Table(s)
    \item 9 Equation(s)
    \item Worked Example(s)
\end{itemize}

\end{tcolorbox}




% =============================================================================
%                    PART 2: CORE CONTENT
% =============================================================================

%%%%%%%%%%%%%%%%%%%%%%%%%%%%%%%%%%%%%%%%%%%%%%%%%%%%%%%%%%%%%%%%%%%%%%%%%%%%%%%
\section{The Fundamental Question}
\label{sec:S-fundamental}
%%%%%%%%%%%%%%%%%%%%%%%%%%%%%%%%%%%%%%%%%%%%%%%%%%%%%%%%%%%%%%%%%%%%%%%%%%%%%%%

\subsection{The Falsifiability Challenge}

A framework that can explain everything explains nothing. This appendix addresses the Popperian challenge by deriving \textbf{specific, falsifiable predictions} from the $(C, \Psi, C^*, K, Q)$ framework.

\subsection{The Core Problem}

The framework claims that:
\begin{enumerate}
    \item Complementarities ($C$) exist between economic inputs
    \item Context ($\Psi$) determines which complementarities are active
    \item The effective complementarity structure $C^*(\Psi)$ varies systematically
    \item Coherence ($K$) measures alignment with $C^*$
    \item All of this determines observable outcomes ($Q$)
\end{enumerate}

\begin{tcolorbox}[colback=red!5!white, colframe=red!75!black,
    title=The Fundamental Question]
\textbf{The critique}: If any empirical finding can be reinterpreted as ``a complementarity moderated by context,'' the framework is unfalsifiable.
\end{tcolorbox}

\subsection{Our Response}

We derive 10 predictions that are:
\begin{itemize}
    \item \textbf{Specific}: Clear about magnitudes and directions
    \item \textbf{Surprising}: Not obvious without the framework
    \item \textbf{Falsifiable}: We specify what evidence would refute them
    \item \textbf{Testable}: Data exists or can be collected
\end{itemize}

\begin{tcolorbox}[colback=blue!5!white, colframe=blue!75!black,
    title=\textbf{Connection to Appendix FRM: The Exclusion Principle}]
This appendix implements the \textbf{Exclusion Principle} established in Appendix FRM
(LIT-FEHR-METHOD). Each prediction below specifies not just \textit{what} EBF predicts,
but crucially \textit{what it does not predict}---i.e., what evidence would falsify it.

\textbf{The methodological stance:}
\begin{itemize}[noitemsep]
    \item $\gamma = 0$ is the null hypothesis (additivity assumed)
    \item Each prediction claims $\gamma \neq 0$ for specific pairs in specific contexts
    \item If observed effects are additive ($\gamma \approx 0$), the prediction fails
\end{itemize}

This makes EBF \textit{Fehr-compatible}: predictions gain value precisely where they can fail.
\end{tcolorbox}


%%%%%%%%%%%%%%%%%%%%%%%%%%%%%%%%%%%%%%%%%%%%%%%%%%%%%%%%%%%%%%%%%%%%%%%%%%%%%%%
\section{Prediction 1: Cash Transfer Effectiveness and Social Capital}
\label{sec:S-prediction1}
%%%%%%%%%%%%%%%%%%%%%%%%%%%%%%%%%%%%%%%%%%%%%%%%%%%%%%%%%%%%%%%%%%%%%%%%%%%%%%%

\paragraph{The Prediction.}
Unconditional cash transfers (UCTs) will show \textbf{smaller} treatment effects in communities with high pre-existing social capital, because informal insurance mechanisms already address the constraints that cash relaxes.

\paragraph{Framework Derivation.}
The framework posits:
\begin{equation}
    \tau_{UCT} = f(\Psi_S) \quad \text{where } \frac{\partial \tau}{\partial \Psi_S} < 0
\end{equation}
High social capital ($\Psi_S$ high) means informal insurance already exists $\Rightarrow$ cash is redundant.

\paragraph{Falsification Criterion.}
The prediction is \textbf{falsified} if UCT effects are:
\begin{itemize}
    \item Equal across high and low social capital communities
    \item \textit{Larger} in high social capital communities
\end{itemize}

\paragraph{Empirical Test.}
Use GiveDirectly data across Kenyan villages, measure pre-treatment social capital (group membership, trust surveys), estimate heterogeneous treatment effects.


%%%%%%%%%%%%%%%%%%%%%%%%%%%%%%%%%%%%%%%%%%%%%%%%%%%%%%%%%%%%%%%%%%%%%%%%%%%%%%%
\section{Prediction 2: Management Practice Transfer Failure}
\label{sec:S-prediction2}
%%%%%%%%%%%%%%%%%%%%%%%%%%%%%%%%%%%%%%%%%%%%%%%%%%%%%%%%%%%%%%%%%%%%%%%%%%%%%%%

\paragraph{The Prediction.}
Management practices that improve productivity in flexible labor markets will show \textbf{zero or negative} effects when transferred to rigid labor markets, because the practices assume adjustment capacity that doesn't exist.

\paragraph{Framework Derivation.}
\begin{equation}
    \tau_{management} = g(\Psi_F) \quad \text{where } \tau(\Psi_F^{low}) \leq 0 < \tau(\Psi_F^{high})
\end{equation}

\paragraph{Falsification Criterion.}
The prediction is \textbf{falsified} if management practices show \textit{uniform positive effects} regardless of labor market flexibility.


%%%%%%%%%%%%%%%%%%%%%%%%%%%%%%%%%%%%%%%%%%%%%%%%%%%%%%%%%%%%%%%%%%%%%%%%%%%%%%%
\section{Prediction 3: Information Intervention Asymmetry}
\label{sec:S-prediction3}
%%%%%%%%%%%%%%%%%%%%%%%%%%%%%%%%%%%%%%%%%%%%%%%%%%%%%%%%%%%%%%%%%%%%%%%%%%%%%%%

\paragraph{The Prediction.}
Information interventions work only when:
\begin{enumerate}
    \item Pre-existing information quality is \textbf{low} ($\Psi_K$ low) AND
    \item Ability to act on information is \textbf{high} ($\Psi_F$ high)
\end{enumerate}

\paragraph{Framework Derivation.}
\begin{equation}
    \tau_{info} = h(\Psi_K, \Psi_F) \approx (1 - \Psi_K) \cdot \Psi_F
\end{equation}
Information is valuable only when novel \textit{and} actionable.

\paragraph{Falsification Criterion.}
The prediction is \textbf{falsified} if information interventions work equally well:
\begin{itemize}
    \item In high-information and low-information environments
    \item In constrained and unconstrained populations
\end{itemize}


%%%%%%%%%%%%%%%%%%%%%%%%%%%%%%%%%%%%%%%%%%%%%%%%%%%%%%%%%%%%%%%%%%%%%%%%%%%%%%%
\section{Prediction 4: Financial Crisis Political Lag}
\label{sec:S-prediction4}
%%%%%%%%%%%%%%%%%%%%%%%%%%%%%%%%%%%%%%%%%%%%%%%%%%%%%%%%%%%%%%%%%%%%%%%%%%%%%%%

\paragraph{The Prediction.}
Major financial crises will be followed by political crises (populism, institutional instability) with a lag of 5--10 years, because financial shocks propagate through social trust ($\Psi_S$) before reaching political outcomes.

\paragraph{Framework Derivation.}
\begin{align}
    \text{Financial shock at } t &\Rightarrow \Psi_E \downarrow \text{ immediately} \\
    &\Rightarrow \Psi_S \downarrow \text{ with lag } \approx 3\text{--}5 \text{ years} \\
    &\Rightarrow \Psi_I \downarrow \text{ with lag } \approx 5\text{--}10 \text{ years}
\end{align}

\paragraph{Falsification Criterion.}
The prediction is \textbf{falsified} if:
\begin{itemize}
    \item Political crises occur immediately after financial crises (no lag)
    \item Political crises occur with lags $> 15$ years or $< 3$ years
    \item No systematic relationship between financial and political crises
\end{itemize}


%%%%%%%%%%%%%%%%%%%%%%%%%%%%%%%%%%%%%%%%%%%%%%%%%%%%%%%%%%%%%%%%%%%%%%%%%%%%%%%
\section{Prediction 5: Recovery Time Convexity}
\label{sec:S-prediction5}
%%%%%%%%%%%%%%%%%%%%%%%%%%%%%%%%%%%%%%%%%%%%%%%%%%%%%%%%%%%%%%%%%%%%%%%%%%%%%%%

\paragraph{The Prediction.}
Recovery time from crises is \textbf{convex} in crisis depth: a crisis twice as deep takes more than twice as long to recover from.

\paragraph{Framework Derivation.}
Coherence rebuilding is harder when starting from lower $K$:
\begin{equation}
    T_{recovery} = \frac{c}{K_{crisis}} \cdot \log\left(\frac{K^*}{K_{crisis}}\right) \quad \text{(convex in } K_{crisis}\text{)}
\end{equation}

\paragraph{Falsification Criterion.}
The prediction is \textbf{falsified} if recovery time is:
\begin{itemize}
    \item Linear in crisis depth
    \item Concave in crisis depth (deeper crises recover faster per unit)
\end{itemize}


%%%%%%%%%%%%%%%%%%%%%%%%%%%%%%%%%%%%%%%%%%%%%%%%%%%%%%%%%%%%%%%%%%%%%%%%%%%%%%%
\section{Prediction 6: Trade War Hysteresis}
\label{sec:S-prediction6}
%%%%%%%%%%%%%%%%%%%%%%%%%%%%%%%%%%%%%%%%%%%%%%%%%%%%%%%%%%%%%%%%%%%%%%%%%%%%%%%

\paragraph{The Prediction.}
Trade wars create persistent effects even after tariffs are removed, because they damage $\Psi_{trust}$ which recovers slowly.

\paragraph{Framework Derivation.}
\begin{equation}
    \Psi_{trust}(t) = \Psi_{trust}^* - \alpha \cdot \mathbb{1}_{tariff} - \beta \cdot \sum_{\tau < t} \mathbb{1}_{tariff,\tau} \cdot e^{-\lambda(t-\tau)}
\end{equation}
The second term captures hysteresis: past tariffs leave scars.

\paragraph{Falsification Criterion.}
The prediction is \textbf{falsified} if trade volumes return to pre-war levels within 1--2 years of tariff removal.


%%%%%%%%%%%%%%%%%%%%%%%%%%%%%%%%%%%%%%%%%%%%%%%%%%%%%%%%%%%%%%%%%%%%%%%%%%%%%%%
\section{Prediction 7: Technology-Culture Divergence}
\label{sec:S-prediction7}
%%%%%%%%%%%%%%%%%%%%%%%%%%%%%%%%%%%%%%%%%%%%%%%%%%%%%%%%%%%%%%%%%%%%%%%%%%%%%%%

\paragraph{The Prediction.}
Rapid technological change ($\dot{\Psi}_{tech}$ high) creates ``cultural lag''---the gap between technical capability and social norms widens, creating predictable backlash.

\paragraph{Framework Derivation.}
\begin{equation}
    \text{Coherence Gap} = |\Psi_{tech} - \Psi_{norm}| \quad \text{increases when } \dot{\Psi}_{tech} > \dot{\Psi}_{norm}
\end{equation}

\paragraph{Falsification Criterion.}
The prediction is \textbf{falsified} if:
\begin{itemize}
    \item Technology adoption faces no social resistance
    \item Resistance is uncorrelated with pace of change
\end{itemize}


%%%%%%%%%%%%%%%%%%%%%%%%%%%%%%%%%%%%%%%%%%%%%%%%%%%%%%%%%%%%%%%%%%%%%%%%%%%%%%%
\section{Prediction 8: Multi-Level Coherence Conflict}
\label{sec:S-prediction8}
%%%%%%%%%%%%%%%%%%%%%%%%%%%%%%%%%%%%%%%%%%%%%%%%%%%%%%%%%%%%%%%%%%%%%%%%%%%%%%%

\paragraph{The Prediction.}
Policies that optimize coherence at one level (e.g., national) may destroy coherence at another level (e.g., global), creating predictable tensions.

\paragraph{Example.}
Climate policy: High $K_{national}$ (domestic political coherence) may require low $K_{global}$ (defection from international agreements).

\paragraph{Falsification Criterion.}
The prediction is \textbf{falsified} if coherence at different levels is always positively correlated.


%%%%%%%%%%%%%%%%%%%%%%%%%%%%%%%%%%%%%%%%%%%%%%%%%%%%%%%%%%%%%%%%%%%%%%%%%%%%%%%
\section{Prediction 9: Identity-Economics Non-Substitutability}
\label{sec:S-prediction9}
%%%%%%%%%%%%%%%%%%%%%%%%%%%%%%%%%%%%%%%%%%%%%%%%%%%%%%%%%%%%%%%%%%%%%%%%%%%%%%%

\paragraph{The Prediction.}
Economic compensation cannot substitute for identity harms. Offers that threaten identity will be rejected even at significant economic cost.

\paragraph{Framework Derivation.}
In FEPSDE: $U^{IDN}$ and $U^F$ are not substitutes:
\begin{equation}
    \frac{\partial^2 U}{\partial U^{IDN} \partial U^F} < 0 \quad \text{(non-compensating)}
\end{equation}

\paragraph{Falsification Criterion.}
The prediction is \textbf{falsified} if sufficiently large economic compensation induces acceptance of identity-threatening offers.


%%%%%%%%%%%%%%%%%%%%%%%%%%%%%%%%%%%%%%%%%%%%%%%%%%%%%%%%%%%%%%%%%%%%%%%%%%%%%%%
\section{Prediction 10: Coherent Traps}
\label{sec:S-prediction10}
%%%%%%%%%%%%%%%%%%%%%%%%%%%%%%%%%%%%%%%%%%%%%%%%%%%%%%%%%%%%%%%%%%%%%%%%%%%%%%%

\paragraph{The Prediction.}
Systems can be stably trapped in low-quality equilibria ($K \approx 1$ but $Q$ low). Escape requires ``big push'' interventions that temporarily reduce coherence.

\paragraph{Framework Derivation.}
Non-concavity creates multiple local maxima. Moving from one to another requires crossing valleys.

\paragraph{Falsification Criterion.}
The prediction is \textbf{falsified} if:
\begin{itemize}
    \item All stable equilibria are high-quality
    \item Gradual change always works (no need for big push)
\end{itemize}


%%%%%%%%%%%%%%%%%%%%%%%%%%%%%%%%%%%%%%%%%%%%%%%%%%%%%%%%%%%%%%%%%%%%%%%%%%%%%%%
\section{Key Results: The Prediction Matrix}
\label{sec:S-results}
%%%%%%%%%%%%%%%%%%%%%%%%%%%%%%%%%%%%%%%%%%%%%%%%%%%%%%%%%%%%%%%%%%%%%%%%%%%%%%%

\begin{table}[htbp]
\centering
\caption{Summary of Ten Falsifiable Predictions}
\label{tab:S-predictions}
\renewcommand{\arraystretch}{1.2}
\begin{tabular}{@{}clll@{}}
\toprule
\textbf{\#} & \textbf{Prediction} & \textbf{Key $\Psi$} & \textbf{Falsification} \\
\midrule
1 & UCT $\times$ Social Capital & $\Psi_S$ & $\tau$ equal across $\Psi_S$ \\
2 & Management transfer & $\Psi_F$ & $\tau > 0$ in rigid markets \\
3 & Information asymmetry & $\Psi_K, \Psi_F$ & $\tau$ equal across $\Psi_K$ \\
4 & Financial-political lag & $\Psi_S \to \Psi_I$ & No lag or immediate \\
5 & Recovery convexity & $K$ & Linear recovery \\
6 & Trade war hysteresis & $\Psi_{trust}$ & Immediate reversal \\
7 & Tech-culture divergence & $\dot{\Psi}$ gap & No backlash \\
8 & Multi-level conflict & $K_{level}$ & Always aligned \\
9 & Identity non-substitution & $U^{IDN}$ & Compensation works \\
10 & Coherent traps & $K, Q$ & No stable low-$Q$ \\
\bottomrule
\end{tabular}
\end{table}

These 10 predictions are specific, falsifiable, and distinguish the framework from alternatives. Each generates testable hypotheses that could refute the framework if disconfirmed.


%%%%%%%%%%%%%%%%%%%%%%%%%%%%%%%%%%%%%%%%%%%%%%%%%%%%%%%%%%%%%%%%%%%%%%%%%%%%%%%
\section{Integration with Other Appendices}
\label{sec:S-integration}
%%%%%%%%%%%%%%%%%%%%%%%%%%%%%%%%%%%%%%%%%%%%%%%%%%%%%%%%%%%%%%%%%%%%%%%%%%%%%%%

\subsection{Connection to Evaluation Protocol (Appendix THL1)}

\begin{tcolorbox}[colback=yellow!5!white, colframe=yellow!75!black,
    title=Integration: Appendix SUN $\leftrightarrow$ Appendix THL1]
\textbf{What this appendix provides:}
\begin{itemize}[nosep]
\item Specific testable hypotheses for empirical validation
\item Falsification criteria for each prediction
\item Framework-derived mathematical formulations
\end{itemize}

\textbf{What it requires:}
\begin{itemize}[nosep]
\item Appendix THL1: Evaluation methodology and data collection protocols
\item Appendix EPS: Protocol for learning from failed predictions
\end{itemize}
\end{tcolorbox}

\subsection{Connection to Epistemics of Deviation (Appendix EPS)}

When predictions fail, Appendix EPS provides the diagnostic protocol for systematic learning from those failures. This creates a feedback loop: failed predictions become inputs for framework refinement.


% =============================================================================
%                    PART 3: APPLICATION
% =============================================================================

%%%%%%%%%%%%%%%%%%%%%%%%%%%%%%%%%%%%%%%%%%%%%%%%%%%%%%%%%%%%%%%%%%%%%%%%%%%%%%%
\section{Worked Example}
\label{sec:S-example}
%%%%%%%%%%%%%%%%%%%%%%%%%%%%%%%%%%%%%%%%%%%%%%%%%%%%%%%%%%%%%%%%%%%%%%%%%%%%%%%

\subsection{Testing Prediction 1: Cash Transfers and Social Capital}

\paragraph{Setup.}
A researcher wants to test whether UCT effects vary with pre-existing social capital using GiveDirectly data from Kenyan villages.

\paragraph{Application.}
\begin{enumerate}
\item \textbf{Measure pre-treatment social capital} ($\Psi_S$):
    \begin{itemize}
    \item Group membership density
    \item Trust survey responses
    \item Informal lending network metrics
    \end{itemize}
\item \textbf{Estimate heterogeneous treatment effects}:
    \begin{equation}
    Y_i = \alpha + \tau \cdot T_i + \beta \cdot \Psi_{S,i} + \gamma \cdot (T_i \times \Psi_{S,i}) + \epsilon_i
    \end{equation}
\item \textbf{Test prediction}: The framework predicts $\gamma < 0$ (interaction is negative)
\item \textbf{Falsification check}: If $\gamma \geq 0$ or $\gamma$ is not significantly different from zero, prediction is falsified
\end{enumerate}

\paragraph{Result.}
If $\hat{\gamma} = -0.15$ with $p < 0.05$, the prediction is supported. If $\hat{\gamma} = 0.02$ with $p > 0.10$, the prediction is falsified.

\paragraph{Discussion.}
This example demonstrates how each prediction can be operationalized with specific data, estimation strategies, and decision rules. The framework commits to specific empirical content, making it genuinely testable.


%%%%%%%%%%%%%%%%%%%%%%%%%%%%%%%%%%%%%%%%%%%%%%%%%%%%%%%%%%%%%%%%%%%%%%%%%%%%%%%
\section{Implications for Practice}
\label{sec:S-practice}
%%%%%%%%%%%%%%%%%%%%%%%%%%%%%%%%%%%%%%%%%%%%%%%%%%%%%%%%%%%%%%%%%%%%%%%%%%%%%%%

\begin{tcolorbox}[colback=green!5!white, colframe=green!75!black,
    title=Practical Implications]
\begin{enumerate}
\item \textbf{Policy Design}: Predictions 1-3 suggest interventions must be tailored to context---one-size-fits-all approaches will fail
\item \textbf{Crisis Management}: Predictions 4-6 imply long planning horizons are needed; crises have delayed and persistent effects
\item \textbf{Technology Adoption}: Prediction 7 suggests change management must address cultural adaptation, not just technical implementation
\item \textbf{Compensation Design}: Prediction 9 warns that economic incentives cannot solve identity-related conflicts
\item \textbf{Development Strategy}: Prediction 10 suggests incremental reform may be insufficient---big push may be necessary
\end{enumerate}
\end{tcolorbox}


% =============================================================================
%                    PART 4: BACK MATTER
% =============================================================================

%%%%%%%%%%%%%%%%%%%%%%%%%%%%%%%%%%%%%%%%%%%%%%%%%%%%%%%%%%%%%%%%%%%%%%%%%%%%%%%
\section{Summary}
\label{sec:S-summary}
%%%%%%%%%%%%%%%%%%%%%%%%%%%%%%%%%%%%%%%%%%%%%%%%%%%%%%%%%%%%%%%%%%%%%%%%%%%%%%%

\begin{tcolorbox}[colback=green!10!white, colframe=green!75!black,
    title=\textbf{Appendix SUN: Summary}]

\textbf{Core Question:} What specific predictions does EBF make that could falsify the framework?

\textbf{Main Contribution:}
\begin{itemize}[nosep]
    \item 10 specific, falsifiable predictions derived from $(C, \Psi, C^*, K, Q)$
    \item Each prediction specifies direction, magnitude, and falsification criterion
    \item Predictions span policy, macro, technology, identity, and development domains
\end{itemize}

\textbf{Key Outputs:}
\begin{itemize}[nosep]
    \item Prediction matrix: Summary table with key variables and falsification criteria
    \item Mathematical formulations: Testable equations for each prediction
    \item Empirical test designs: Specific data and methods for validation
\end{itemize}

\textbf{Integration:} This appendix connects to Appendix THL1 (testing methodology), Appendix CAM (metatheory), and Appendix EPS (learning from failures).

\end{tcolorbox}


%%%%%%%%%%%%%%%%%%%%%%%%%%%%%%%%%%%%%%%%%%%%%%%%%%%%%%%%%%%%%%%%%%%%%%%%%%%%%%%
\section{Glossary of Symbols}
\label{sec:S-glossary}
%%%%%%%%%%%%%%%%%%%%%%%%%%%%%%%%%%%%%%%%%%%%%%%%%%%%%%%%%%%%%%%%%%%%%%%%%%%%%%%

\begin{tcolorbox}[colback=blue!5!white, colframe=blue!75!black]
\textbf{Central Reference:} For complete symbol definitions, see
\textbf{Appendix GLS} (\textit{Glossary and Notation}, \S\ref{app:glossary}).

This section lists only symbols introduced or specialized in this appendix.
\end{tcolorbox}

\vspace{0.5em}

\begin{table}[htbp]
\centering
\caption{Symbols in Appendix SUN}
\label{tab:S-symbols}
\renewcommand{\arraystretch}{1.3}
\begin{tabular}{@{}clp{6cm}c@{}}
\toprule
\textbf{Symbol} & \textbf{Name} & \textbf{Definition} & \textbf{Also in G?} \\
\midrule
$C$ & Complementarity & Structural complementarity between inputs & \checkmark \\
$\Psi$ & Context Vector & 8-dimensional context state & \checkmark \\
$C^*(\Psi)$ & Effective Complementarity & Context-activated complementarity structure & \checkmark \\
$K$ & Coherence & Alignment with effective complementarity & \checkmark \\
$Q$ & Outcomes & Observable outcome variable & \checkmark \\
$\tau$ & Treatment Effect & Causal effect of intervention & NEW \\
$\Psi_S$ & Social Capital & Social context dimension & \checkmark \\
$\Psi_F$ & Flexibility & Market/institutional flexibility dimension & \checkmark \\
$\Psi_K$ & Knowledge & Information quality dimension & \checkmark \\
$\Psi_E$ & Economic & Economic context dimension & \checkmark \\
$\Psi_I$ & Institutional & Institutional context dimension & \checkmark \\
$U^{IDN}$ & Identity Utility & Identity component of utility & \checkmark \\
$U^F$ & Financial Utility & Financial component of utility & \checkmark \\
$T_{recovery}$ & Recovery Time & Time to return to equilibrium & NEW \\
\bottomrule
\end{tabular}
\end{table}


%%%%%%%%%%%%%%%%%%%%%%%%%%%%%%%%%%%%%%%%%%%%%%%%%%%%%%%%%%%%%%%%%%%%%%%%%%%%%%%
\section{Critical Foundations and Objections}
\label{sec:S-foundations}
%%%%%%%%%%%%%%%%%%%%%%%%%%%%%%%%%%%%%%%%%%%%%%%%%%%%%%%%%%%%%%%%%%%%%%%%%%%%%%%

\subsection{Objection 1: Post-Hoc Rationalization}

\textbf{Objection:} These predictions were derived after observing the phenomena---they are post-hoc rationalizations, not genuine predictions.

\textbf{Response:} Several predictions (e.g., recovery convexity, trade war hysteresis) generate specific quantitative relationships that were not used in framework construction. The framework commits to magnitudes and directions that can be tested against future data. Moreover, the predictions are stated with specific falsification criteria---if evidence contradicts them, the framework must be revised.

\subsection{Objection 2: Too Many Free Parameters}

\textbf{Objection:} With context variables $\Psi$, complementarities $C$, and coherence $K$, any result can be explained by adjusting parameters.

\textbf{Response:} Each prediction constrains the parameters. For example, Prediction 1 requires $\frac{\partial \tau}{\partial \Psi_S} < 0$---the framework cannot accommodate $\frac{\partial \tau}{\partial \Psi_S} > 0$ without revision. The predictions collectively constrain the parameter space significantly.

\subsection{Objection 3: Domain-Specific Not General}

\textbf{Objection:} These predictions are domain-specific (development economics, macroeconomics, etc.) and don't test the general framework.

\textbf{Response:} The predictions are derived from general framework components ($C$, $\Psi$, $K$, $Q$) applied to specific domains. If predictions fail across multiple domains, this falsifies the general framework. Domain-specificity is necessary for testability---abstract claims cannot be tested.


%%%%%%%%%%%%%%%%%%%%%%%%%%%%%%%%%%%%%%%%%%%%%%%%%%%%%%%%%%%%%%%%%%%%%%%%%%%%%%%
\section{Open Issues and Research Agenda}
\label{sec:S-open}
%%%%%%%%%%%%%%%%%%%%%%%%%%%%%%%%%%%%%%%%%%%%%%%%%%%%%%%%%%%%%%%%%%%%%%%%%%%%%%%

\subsection{Methodological Priorities}

\begin{enumerate}
    \item \textbf{Operationalization}: Develop standardized measures for $\Psi$ dimensions across contexts
    \item \textbf{Power Analysis}: Determine sample sizes needed to detect predicted effect sizes
    \item \textbf{Meta-Analysis}: Compile existing evidence that tests these predictions
\end{enumerate}

\subsection{Empirical Priorities}

\begin{enumerate}
    \item \textbf{Prediction 1 Test}: Partner with GiveDirectly for heterogeneity analysis
    \item \textbf{Prediction 4 Test}: Historical analysis of financial-political crisis timing (2008, 1997 Asian, 1929)
    \item \textbf{Prediction 9 Test}: Experimental tests of identity-economics trade-offs
    \item \textbf{Cross-Prediction}: Test whether predictions hold across cultural contexts
\end{enumerate}


%%%%%%%%%%%%%%%%%%%%%%%%%%%%%%%%%%%%%%%%%%%%%%%%%%%%%%%%%%%%%%%%%%%%%%%%%%%%%%%
\section{References}

% Link to master bibliography
\nocite{bcm_master}

\label{sec:S-references}
%%%%%%%%%%%%%%%%%%%%%%%%%%%%%%%%%%%%%%%%%%%%%%%%%%%%%%%%%%%%%%%%%%%%%%%%%%%%%%%

\begin{tcolorbox}[colback=blue!5!white, colframe=blue!75!black]
\textbf{Master Bibliography:} For complete reference details, see
\texttt{ebf\_master\_references.bib} in the repository.

This section lists appendix-specific references and data sources.
\end{tcolorbox}

\vspace{0.5em}

\subsection{Key References for This Appendix}

The following works are particularly relevant for the content of this appendix:

\begin{itemize}[nosep]
    \item \textbf{Falsifiability}: Popper (1959), Lakatos (1970)
    \item \textbf{Cash Transfers}: Haushofer \& Shapiro (2016), Banerjee et al. (2015)
    \item \textbf{Management Practices}: Bloom et al. (2013), Bloom \& Van Reenen (2007)
    \item \textbf{Financial Crises}: Funke et al. (2016), Reinhart \& Rogoff (2009)
    \item \textbf{Identity Economics}: Akerlof \& Kranton (2000, 2010)
    \item \textbf{Big Push}: Murphy, Shleifer \& Vishny (1989), Rosenstein-Rodan (1943)
\end{itemize}

\subsection{Appendix-Specific References}

\begin{thebibliography}{99}

\bibitem{popper1959} Popper, K. (1959). The Logic of Scientific Discovery. Hutchinson.

\bibitem{lakatos1970} Lakatos, I. (1970). Falsification and the Methodology of Scientific Research Programmes. In Lakatos \& Musgrave (Eds.), Criticism and the Growth of Knowledge.

\bibitem{haushofer2016} Haushofer, J., \& Shapiro, J. (2016). The Short-Term Impact of Unconditional Cash Transfers to the Poor. \textit{Quarterly Journal of Economics}, 131(4), 1973--2042.

\bibitem{bloom2013} Bloom, N., Eifert, B., Mahajan, A., McKenzie, D., \& Roberts, J. (2013). Does Management Matter? Evidence from India. \textit{Quarterly Journal of Economics}, 128(1), 1--51.

\bibitem{funke2016} Funke, M., Schularick, M., \& Trebesch, C. (2016). Going to Extremes: Politics after Financial Crises, 1870--2014. \textit{European Economic Review}, 88, 227--260.

\bibitem{akerlof2000} Akerlof, G. A., \& Kranton, R. E. (2000). Economics and Identity. \textit{Quarterly Journal of Economics}, 115(3), 715--753.

\bibitem{murphy1989} Murphy, K. M., Shleifer, A., \& Vishny, R. W. (1989). Industrialization and the Big Push. \textit{Journal of Political Economy}, 97(5), 1003--1026.

\end{thebibliography}


%%%%%%%%%%%%%%%%%%%%%%%%%%%%%%%%%%%%%%%%%%%%%%%%%%%%%%%%%%%%%%%%%%%%%%%%%%%%%%%
% CROSS-REFERENCE FOOTER
%%%%%%%%%%%%%%%%%%%%%%%%%%%%%%%%%%%%%%%%%%%%%%%%%%%%%%%%%%%%%%%%%%%%%%%%%%%%%%%

\vspace{1em}
\hrule
\vspace{0.5em}

\noindent\textit{Cross-references:}
Appendix GLS (Central Glossary),
Appendix THL1 (Evaluation Protocol),
Appendix CAM (Metatheory),
Appendix EPS (Epistemics of Deviation---when predictions fail, Appendix EPS provides the diagnostic protocol for systematic learning from those failures),
Appendix HOW (CORE-HOW),
Appendix CTW (CORE-WHEN),
Appendix WAT (CORE-WHAT).

\noindent\textit{Master Bibliography:} \texttt{ebf\_master\_references.bib}

% =============================================================================
% END OF APPENDIX S
% =============================================================================
